Bitcoin trades continuously and is characterized by pronounced
and rapidly changing volatility. Reliable short-horizon
volatility forecasts are relevant to risk management, position
sizing, derivatives pricing, and intraday trading decisions.
Previous studies show that high-frequency realized-volatility
measures can support Bitcoin volatility forecasting, while
jumps and structural breaks may materially affect predictive
performance \citep{bergsli2022,shen2020}. These features
motivate models that capture persistent volatility dynamics
and allow forecasting relationships to evolve over time.

High-frequency returns provide a direct basis for measuring
short-term price variation. Realized variance aggregates
squared intraday returns and approximates the variation
accumulated over a specified interval \citep{andersen2003}.
This study constructs hourly realized variance from
five-minute returns and uses its logarithm as the forecasting
target. The resulting measure retains substantially more
intraday information than volatility proxies based on a single
hourly or daily return.

Autoregressive models provide a natural benchmark because
volatility is strongly persistent. However, a fixed collection
of individual lags may represent multi-scale dependence
inefficiently. The heterogeneous autoregressive model addresses
this issue by combining realized-volatility components measured
over different horizons \citep{corsi2009}. The specification
used here adapts that framework to hourly data through
one-hour, six-hour, daily, and weekly volatility components.

The analysis also considers whether current market activity
contains predictive information beyond lagged volatility. The
mixture-of-distributions hypothesis links trading intensity and
price variability through the arrival of information
\citep{clark1973}, while empirical evidence documents a close
relationship between transaction counts and volatility
\citep{jones1994}. Three additional predictors are therefore
included: downside realized semivariance, the number of trades,
and an aggressive-buying proxy. Downside realized
semivariance isolates the contribution of negative intraday
returns to total realized variation
\citep{barndorffnielsen2010}.

A fixed-parameter model imposes a constant relationship between
these predictors and future volatility. Such an assumption may
be restrictive in cryptocurrency markets, where participation,
liquidity, and volatility conditions can change substantially
over time \citep{shen2020}. To accommodate gradual parameter
variation, this study estimates a time-varying parameter HAR-X
model in state-space form. The coefficients follow random-walk
state equations and are updated sequentially through the
Kalman filter \citep{kalman1960,durbin2012}.

The empirical analysis uses five-minute BTCUSDT
USD-margined perpetual futures data from January 2022 to
December 2025. Four forecasting models are compared: a
persistence benchmark, AR(24), Fixed HAR-X, and TVP-HAR-X.
The sample is divided chronologically into training,
validation, and final test periods. Preprocessing parameters
and model specifications are determined using the training and
validation samples, and the 2025 test sample is reserved for a
single final evaluation.

The study addresses two questions. First, do multi-horizon
volatility components and high-frequency market activity
improve one-hour-ahead forecasts relative to persistence and
AR(24)? Second, does coefficient variation provide additional
predictive value beyond a fixed-parameter HAR-X model using the
same information set?

The results show that both HAR-X specifications substantially
outperform persistence and AR(24). Fixed HAR-X achieves the
lowest test QLIKE, while TVP-HAR-X records the lowest MAE and
RMSE. Trading activity contributes the clearest incremental
information among the additional predictors. Coefficient
variation improves average forecast accuracy on the log scale,
although it does not improve the primary variance-scale loss.

This study makes three contributions. It evaluates the proposed
models against both persistence and a strong autoregressive
benchmark; it combines multi-horizon realized volatility with
hourly market-activity variables; and it compares fixed and
time-varying coefficients using an identical predictor set.
The last comparison isolates the forecasting value of parameter
adaptation from differences in observable information.

The remainder of the report is organized as follows.
Section~\ref{sec:data} describes the data and variable
construction. Section~\ref{sec:forecasting-design} presents the
forecasting design and evaluation criteria.
Section~\ref{sec:methodology} introduces the four models.
Section~\ref{sec:results} reports the empirical results, and
Section~\ref{sec:discussion} discusses their interpretation
and limitations. Section~\ref{sec:conclusion} concludes.